
This study provides the first quantitative measurement of overlap between static and execution feedback signals for LLM code generation. The primary findings are:

\begin{enumerate}
\item \textbf{Complete orthogonality:} Jaccard similarity = 0.0 between static-detected and execution-detected error sets, indicating these feedback types detect entirely different error classes.
\end{enumerate}

\begin{enumerate}
\item \textbf{High static coverage:} Static analysis detects structural errors in 98.6\% of test-failing code.
\end{enumerate}

\begin{enumerate}
\item \textbf{Behavioral detection untested:} The use of canonical solutions prevented validation of the hypothesis that execution detects behavioral errors in static-clean code (observed rate: 0.3\% vs. 40\% threshold).
\end{enumerate}

The complete separation between error classes supports the theoretical basis for combining static and execution feedback in iterative refinement systems. However, empirical validation of combined feedback improvement requires experiments with actual LLM-generated code.

\textbf{Future Work:}
\begin{itemize}
\item Replicate analysis on actual LLM-generated code to test behavioral detection hypothesis
\item Measure combined feedback improvement magnitude
\item Extend to multiple LLM architectures and programming languages
\end{itemize}

